\documentclass{report}
\usepackage{amsmath,amssymb}
\setlength{\parindent}{0mm}
\setlength{\parskip}{1em}
\begin{document}
\begin{center}
\rule{6in}{1pt} \
{\large Van Emden Henson \\
{\bf Tools for spectral near-bipartite community detection}}

Lawrence Livermore National Laboratory \\ Center for Applied Scientific Computing \\ 7000 East Avenue \\ Livermore \\ Ca 94550
\\
{\tt henson5@llnl.gov}\\
Geoffrey Sanders\end{center}

Near-bipartite communities (NBPCs) are an example of dense graph
structure with interesting internal structure. NBPCs are common in
networks involving competition amongst subsets of entities (e.g. a dating
network). We discuss uses of spectral embedding for NBPC detection.
Typical spectral algorithms for detecting traditional community structure
involve the eigenvectors of the graph Laplacian associated with the
smallest non-zero eigenvalues (smoothest eigenmodes). For NBPCs, we study
algorithms based on embeddings using eigenvectors associated with the
largest eigenvalues of the normalized graph Laplacian (most oscillatory
eigenmodes). We present important preprocessing steps, how the
eigenvectors indicate the presence of NBPCs, and post-processing steps
for robust detection. Additionally, we present preliminary convergence
criteria and analysis of iterative eigensolvers for calculating
oscillatory Laplacian eigenmodes. We use several graph models to
illustrate concepts and demonstrate performance of our algorithms. We
include some tests on real-world complex networks.


\end{document}
